% ============================================
% Supplementary Note 3: Simulation Robustness
% ============================================
\subsection*{Supplementary Note 3: Simulation Robustness}
\label{sec:suppl:note3}

\rev{In the main simulations, we use $\mu$ as a minimal one-parameter control of effective interaction strength. For the baseline uniform distribution $U[0,2\mu]$, increasing $\mu$ both moves the average coefficient away from the no-interaction limit at zero and broadens the sampled range of pairwise coefficients. This simple parameter therefore captures two aspects of stronger effective interactions: greater average inhibitory effect and greater heterogeneity among species pairs.}

\rev{We first tested whether the qualitative transition depended on the shape of the coefficient distribution. We compared three distributions with matched mean interaction strength $\mu$: uniform ranging from 0 to $2\mu$ (primary analysis), Gaussian with mean $\mu$ and standard deviation $\mu/\sqrt{3}$ (truncated at 0), and Gamma with mean $\mu$ and shape parameter 3. All three distributions were parameterized to have the same mean $\mu$ and matched variance $\mu^2/3$ (for the uniform distribution $U[0, 2\mu]$, $\text{var} = (2\mu)^2/12 = \mu^2/3$). The qualitative transition from Mixture-dominated to Dominance/Restructuring outcomes with increasing interaction strength was robust across all tested distributions (Supplementary Fig.~4).}

We also tested whether community size (number of species per parental community) affects the frequency of Dominance outcomes. Simulations were performed with parental community sizes ranging from 4 to 48 species per community (\rev{Extended Data Fig.~4}). The qualitative patterns were consistent across community sizes: Dominance frequency remained relatively stable ($\sim$14--26\% at $\mu=0.3$, $\sim$56--67\% at $\mu=0.6$, $\sim$69--78\% at $\mu=0.8$), while Mixture decreased and Restructuring increased with larger species pools. This pattern reflects increased opportunity for competitive exclusion cascades in larger communities, but the qualitative prevalence of Dominance at moderate-to-high interaction strengths is robust to community size.

\rev{Because the baseline gLV model is restricted to non-negative, growth-inhibiting coefficients, we also tested two extensions of the interaction structure while applying the same assembly, coalescence, and outcome-classification pipeline. First, we introduced a mixed-sign facilitative-tail model in which off-diagonal coefficients were sampled from $\alpha_{ij}\sim U[-f\mu,(2+f)\mu]$, preserving $E[\alpha_{ij}]=\mu$ while allowing a small fraction of coefficients to be negative under our sign convention (positive ecological interactions). To prevent runaway growth in the presence of facilitation, we added density-dependent self-limitation, $\dot n_i=n_i(1-(A n)_i-0.1n_i^2)$. Across $f \in \{0,0.10,0.20,0.40,0.80\}$ and $\mu \in \{0.30,0.60,0.80\}$, Dominance still increased with $\mu$, although increasing the facilitative tail shifted some low-$\mu$ outcomes from Mixture toward Dominance and Restructuring (Supplementary Fig.~40).}

\rev{Second, we varied the reciprocal sign/correlation structure between coefficients $\alpha_{ij}$ and $\alpha_{ji}$. For a fraction $|p|$ of unordered species pairs, coefficients were constrained to be antisymmetric exploitation when $p<0$ ($\alpha_{ji}=-\alpha_{ij}$) or symmetric competition when $p>0$ ($\alpha_{ji}=\alpha_{ij}$); remaining pairs were drawn iid from the baseline competitive distribution. This pair-coupling sweep tests whether the Dominance increase depends on independent directed interaction coefficients. The qualitative $\mu$-dependent increase in Dominance persisted across the reciprocal-correlation sweep and in a finer sweep from $p=-1$ to $p=+1$ (Supplementary Fig.~41).}

\rev{Third, to address the specific mutualism question, we converted a controlled fraction of unordered species pairs into weak bidirectional mutualisms. For $q \in \{0,0.10,0.20,0.30,0.40\}$ of species pairs, both reciprocal coefficients were sampled from $\alpha_{ij},\alpha_{ji}\sim U[-0.2\mu,0]$, while all remaining off-diagonal coefficients were sampled from the baseline competitive distribution $U[0,2\mu]$. We used the same density-dependent self-limitation and the same assembly, coalescence, and classification pipeline as in the mixed-sign facilitative-tail analysis. Weak mutualistic pairs did not eliminate the interaction-strength-dependent Dominance trend: at $\mu=0.60$, Dominance was 58.5\%, 70.2\%, 69.8\%, 67.3\%, and 58.7\% for $q=0,0.10,0.20,0.30,0.40$, respectively, and at $\mu=0.80$ it remained 69.4--77.5\% across the same sweep (Supplementary Fig.~42).}

\rev{Fourth, to test which components of this one-parameter interaction-strength axis are required, we separated coefficient mean and coefficient spread in a mean-vs-variance sweep. Off-diagonal coefficients were sampled from $\alpha_{ij}\sim U[m-h,m+h]$ for $m \in \{0,0.2,0.4,0.6,0.8\}$ and $h \in \{0,0.2,0.4,0.6,0.8\}$, giving $\mathrm{std}(\alpha_{ij})=h/\sqrt{3}$. Negative coefficients, corresponding to positive ecological interactions under our sign convention, occur when $h>m$. When $h=0$, all outcomes were Mixture even at $m=0.8$, showing that coefficient mean alone is not sufficient to generate Dominance. Increasing $h$ at fixed $m$ increased Dominance and same-parent concordance, while the strongest Dominance values occurred when both the coefficient mean and coefficient spread were high; across the 25 cells, Dominance correlated more strongly with $h$ (and therefore coefficient standard deviation) than with $m$ (descriptive Pearson $r=0.78$ versus $0.53$; Supplementary Fig.~43). Thus, both average inhibitory effect and coefficient heterogeneity contribute to the effective interaction-strength axis captured by the baseline $\mu$ parameter. Together, these robustness checks indicate that the main simulation result is not restricted to iid non-negative coefficients. Strongly mutualistic or facilitation-dominated systems remain outside the baseline model scope.}
